\documentclass[12pt,a4paper]{article}
\usepackage[utf8]{inputenc}
\usepackage[T5]{fontenc}
\usepackage[vietnamese]{babel}
\usepackage[margin=1.3cm,top=0.8cm,bottom=1.0cm,includehead,includefoot,headheight=18pt,headsep=6pt,footskip=16pt]{geometry}
\usepackage{amsmath,amssymb,array,tabularx,enumitem,fancyhdr,tikz,xcolor,needspace}
\usetikzlibrary{calc,angles,quotes,arrows.meta,babel,patterns}
\definecolor{ink}{RGB}{30,55,70}

\newcommand{\headerQuote}{%
  \ifcase\value{page}\relax
    \or Math gives us hope that every problem has a solution% Trang 1
    \or Math + You = Champion% Trang 2
    \or Where Frustration Becomes Inspiration% Trang 3
    \or Turning Confusion into Clarity% Trang 4
    \or Don't be an X, be a Y of mathematics% Trang 5
    \or Don't get mad. Learn to add% Trang 6
    \or The only way to learn mathematics is to do mathematics% Trang 7
    \or Small steps every day lead to big results% Trang 8
    \or Every expert was once a beginner% Trang 9
    \or Success is the sum of small efforts repeated daily% Trang 10
    \or Mathematics is the music of reason% Trang 11
    \or Believe you can and you're halfway there% Trang 12
    \else Math gives us hope that every problem has a solution%
  \fi
}
\pagestyle{fancy}\fancyhf{}
\lhead{\small\textit{\headerQuote}}
\rhead{\small\textit{Người biên soạn: Dương Gia Bảo}}
\rfoot{\small Trang \thepage/5}
\renewcommand{\headrulewidth}{0.4pt}
\fancypagestyle{dapan}{%
  \fancyhf{}%
  \lhead{\small\textit{\headerQuote}}%
  \rhead{\small\textit{Người biên soạn: Dương Gia Bảo}}%
  \rfoot{\small Trang \the\numexpr\value{page}-5\relax/2}%
  \renewcommand{\headrulewidth}{0.4pt}%
}
\setlength{\headheight}{18pt}
\setlength{\parindent}{0pt}
\setlength{\parskip}{0pt}
\setlength{\tabcolsep}{2pt}
\renewcommand{\arraystretch}{1.05}
\newcommand{\parttitle}[1]{\par\vspace{2.5pt}\needspace{2\baselineskip}{\large\bfseries\color{ink}#1}\par\nopagebreak\vspace{1.5pt}}
\newcommand{\answerbox}{\hfill\textit{Đáp số:}\ \fbox{\rule{0pt}{9pt}\hspace{2.2cm}}}
\setlist[enumerate,1]{label=\textbf{Câu \arabic*.},leftmargin=1.1cm,labelsep=.15cm,itemsep=2.5pt,topsep=1pt,parsep=0pt}
\setlist[enumerate,2]{label=\alph*),leftmargin=.55cm,itemsep=1.5pt,topsep=0.2pt,parsep=0pt}

\begin{document}
% ==========================================
% TRANG 1: TIÊU ĐỀ + CHỦ ĐỀ 1 (CÂU 1 - 8)
% ==========================================
\noindent
\begin{minipage}[c]{0.38\textwidth}
    \centering
    \textbf{LÀM TOÁN ĐỂ HẠNH PHÚC}\\[-2pt]
    \rule{3.2cm}{0.6pt}\\[2pt]
    \textbf{ĐỀ CƯƠNG ÔN TẬP}
\end{minipage}\hfill
\begin{minipage}[c]{0.38\textwidth}
    \centering
    \textbf{ĐỀ CƯƠNG ÔN TẬP GIỮA KÌ I}\\
    \textbf{NĂM HỌC 2026 -- 2027}\\[1pt]
    \textbf{CHUYÊN ĐỀ: 40 CÂU TRẢ LỜI NGẮN}\\
    \textbf{Môn: TOÁN -- Khối lớp 10}\\
    \textit{\small (Mức độ: Vận dụng -- Vận dụng cao)}
\end{minipage}

\vspace{0.1cm}
\noindent
\begin{minipage}[b]{0.78\textwidth}
    \textbf{Họ và tên:}\dotfill\textbf{SBD:}\makebox[3.5cm]{\dotfill}
\end{minipage}\hfill
\begin{minipage}[b]{0.20\textwidth}
    \hfill
    \setlength{\fboxrule}{0.8pt}%
    \fbox{\begin{minipage}{2.3cm}
        \centering
        \textbf{\small Tài liệu}\\
        \textbf{LƯU HÀNH}
    \end{minipage}}
\end{minipage}

\vspace{0.12cm}
\textit{\textbf{Lưu ý:} Học sinh trả lời từ Câu 1 đến Câu 40. Điền kết quả (dưới dạng số nguyên, phân số tối giản hoặc số thập phân theo yêu cầu) vào ô đáp số tương ứng.}

\parttitle{CHỦ ĐỀ 1. MỆNH ĐỀ VÀ TẬP HỢP (CÂU 1 -- 13)}

\begin{enumerate}
\item Cho hai tập hợp $A = (-2; 5)$ và $B = (m - 1; m + 3]$ (với $m$ là tham số thực). Có bao nhiêu giá trị nguyên của $m$ để tập hợp giao $A \cap B$ là một khoảng (hoặc nửa khoảng) có độ dài bằng $3$?
\answerbox

\item Cho hai tập hợp $A = [m - 1; m + 4]$ và $B = (-3; 6)$ (với $m$ là tham số thực). Có bao nhiêu giá trị nguyên của tham số $m$ để $A \subset B$?
\answerbox

\item Cho hai tập hợp $A = (-\infty; m + 2]$ và $B = [2m - 1; +\infty)$ (với $m$ là tham số thực). Tìm giá trị nguyên nhỏ nhất của tham số $m$ để $A \cap B = \varnothing$.
\answerbox

\item Cho tập hợp $A = \{x \in \mathbb{R} \mid x^2 - (2m + 1)x + m^2 + m = 0\}$ và tập hợp $B = (1; 5)$ (với $m$ là tham số thực). Có bao nhiêu giá trị nguyên của $m$ để $A \subset B$?
\answerbox

\item Trong đợt khảo sát $100$ học sinh khối 10 tham gia ba câu lạc bộ Toán học, Tin học và Tiếng Anh, kết quả thu được: có $45$ học sinh tham gia câu lạc bộ Toán học, $40$ học sinh tham gia Tin học, $50$ học sinh tham gia Tiếng Anh; có $20$ học sinh tham gia cả Toán và Tin, $18$ học sinh tham gia cả Tin và Tiếng Anh, $22$ học sinh tham gia cả Toán và Tiếng Anh; có đúng $10$ học sinh tham gia cả ba câu lạc bộ. Hỏi có bao nhiêu học sinh không tham gia bất kỳ câu lạc bộ nào trong ba câu lạc bộ trên?
\answerbox

\item Một trường THPT tổ chức khảo sát năng khiếu thể thao của $120$ học sinh. Kết quả cho thấy: có $70$ học sinh giỏi môn Bóng đá, $65$ học sinh giỏi Bóng chuyền và $60$ học sinh giỏi Cầu lông. Biết rằng mỗi học sinh đều giỏi ít nhất một môn thể thao và có đúng $15$ học sinh giỏi cả ba môn. Hỏi có tất cả bao nhiêu học sinh chỉ giỏi đúng hai môn thể thao?
\answerbox

\item Cho hai tập hợp $A = [m; m + 2]$ và $B = [1; 3] \cup [5; 8]$ (với $m$ là tham số thực). Tính tổng tất cả các giá trị nguyên của tham số $m$ để $A \subset B$.
\answerbox

\item Cho mệnh đề chứa biến $P(m)$: ``$\forall x \in \mathbb{R},\ x^2 - 2(m - 1)x + 2m + 1 > 0$''. Có bao nhiêu giá trị nguyên của tham số $m$ thuộc đoạn $[-10; 10]$ để mệnh đề $P(m)$ là một mệnh đề đúng?
\answerbox
\end{enumerate}

\newpage

% ==========================================
% TRANG 2: CÂU 9 - 13 + CHỦ ĐỀ 2 (CÂU 14 - 16)
% ==========================================
\begin{enumerate}[resume]
\item Cho hai tập hợp $A = (-4; 1]$ và $B = (m - 2; m + 5)$ (với $m$ là tham số thực). Có bao nhiêu giá trị nguyên của tham số $m$ để $A \cap B \ne \varnothing$?
\answerbox

\item Cho hai tập hợp $A = (-3; 4]$ và $B = (-\infty; m) \cup (m + 3; +\infty)$ (với $m$ là tham số thực). Có bao nhiêu giá trị nguyên của tham số $m$ thuộc đoạn $[-10; 10]$ để $A \subset B$?
\answerbox

\item Trong một lớp học có $40$ học sinh, có $30$ học sinh thích môn Toán và $25$ học sinh thích môn Ngữ văn. Gọi $x$ là số học sinh thích cả hai môn Toán và Ngữ văn. Giá trị nhỏ nhất của $x$ bằng bao nhiêu?
\answerbox

\item Cho tập hợp $X = \{1; 2; 3; 4; 5; 6; 7\}$. Có tất cả bao nhiêu tập hợp con $A$ của tập hợp $X$ thoả mãn điều kiện $A$ chứa cả hai phần tử $1, 2$ và không chứa phần tử $7$?
\answerbox

\item Cho hai tập hợp $A = (-\infty; 2m - 3)$ và $B = [m + 1; +\infty)$ (với $m$ là tham số thực). Có bao nhiêu giá trị nguyên của tham số $m$ thuộc đoạn $[-5; 10]$ để $A \cup B = \mathbb{R}$?
\answerbox
\end{enumerate}

\parttitle{CHỦ ĐỀ 2. BẤT PHƯƠNG TRÌNH VÀ HỆ BPT BẬC NHẤT HAI ẨN (CÂU 14 -- 26)}

\begin{enumerate}[resume]
\item Cho miền tứ giác $OABC$ (kể cả biên) có toạ độ các đỉnh là gốc toạ độ $O(0; 0)$, $A(4; 0)$, $B(3; 3)$ và $C(0; 4)$. Tìm giá trị lớn nhất của biểu thức $F(x, y) = 4x + 3y$ với cặp số $(x; y)$ thuộc miền tứ giác $OABC$.
\answerbox

\item Cho cặp số $(x; y)$ thoả mãn hệ bất phương trình:
\[
\begin{cases}
x + 2y \ge 6 \\
2x + y \ge 6 \\
x \ge 0,\ y \ge 0
\end{cases}
\]
Tìm giá trị nhỏ nhất của biểu thức $F(x, y) = 3x + 4y$.
\answerbox

\item Tính diện tích của miền nghiệm xác định bởi hệ bất phương trình bậc nhất hai ẩn:
\[
\begin{cases}
3x + y \le 12 \\
x + 3y \le 12 \\
x \ge 0,\ y \ge 0
\end{cases}
\]
\answerbox
\end{enumerate}

\newpage

% ==========================================
% TRANG 3: CÂU 17 - 24
% ==========================================
\begin{enumerate}[resume]
\item Một xưởng mộc mỹ nghệ sản xuất hai loại sản phẩm: Bàn trà và Ghế tựa. Để sản xuất một bàn trà cần $3$ giờ tiện gỗ và $2$ giờ sơn phủ; để sản xuất một ghế tựa cần $1$ giờ tiện gỗ và $2$ giờ sơn phủ. Mỗi ngày xưởng có tối đa $18$ giờ làm việc của máy tiện và $16$ giờ của máy sơn. Biết mỗi bàn trà mang lại lợi nhuận $400$ nghìn đồng, mỗi ghế tựa mang lại lợi nhuận $200$ nghìn đồng. Hỏi lợi nhuận tối đa trong một ngày của xưởng bằng bao nhiêu nghìn đồng?
\answerbox

\item Một doanh nghiệp cần thuê xe để vận chuyển $140$ công nhân và $9$ tấn hàng. Nơi cho thuê có hai loại xe: loại $A$ (chở được tối đa $20$ người và $0{,}6$ tấn hàng, giá thuê $4$ triệu đồng/xe, có tối đa $8$ xe) và loại $B$ (chở được tối đa $10$ người và $1{,}5$ tấn hàng, giá thuê $3$ triệu đồng/xe, có tối đa $9$ xe). Hỏi chi phí thuê xe nhỏ nhất là bao nhiêu triệu đồng?
\answerbox

\item Có bao nhiêu cặp số nguyên $(x; y)$ thoả mãn hệ bất phương trình:
\[
\begin{cases}
x \ge 0,\ y \ge 0 \\
x + y \le 4 \\
2x - y \ge 0
\end{cases}
\]
\answerbox

\item Một bác sĩ dinh dưỡng tư vấn cho một bệnh nhân cần mua thịt bò và thịt lợn sao cho tổng lượng dinh dưỡng nạp vào đạt ít nhất $2400$ đơn vị protein và ít nhất $1600$ đơn vị chất béo. Mỗi kilôgam thịt bò chứa $600$ đơn vị protein và $200$ đơn vị chất béo, giá $220$ nghìn đồng/kg (khuyến nghị mua không quá $5\text{ kg}$). Mỗi kilôgam thịt lợn chứa $400$ đơn vị protein và $400$ đơn vị chất béo, giá $180$ nghìn đồng/kg (mua không quá $5\text{ kg}$). Hỏi số tiền ít nhất bệnh nhân đó phải chi là bao nhiêu nghìn đồng?
\answerbox

\item Cho miền tam giác $ABC$ (kể cả biên) có toạ độ các đỉnh là $A(1; 1)$, $B(4; 1)$, $C(1; 4)$. Tìm giá trị nhỏ nhất của tham số $m$ để bất phương trình $2x + 3y \le m$ nghiệm đúng với mọi điểm $(x; y)$ thuộc miền tam giác $ABC$.
\answerbox

\item Một hợp tác xã nông nghiệp dự định dùng $10\text{ ha}$ đất để trồng hai loại cây trồng là Lúa và Ngô. Chi phí công lao động để chăm sóc $1\text{ ha}$ lúa là $10$ công, lợi nhuận thu được là $30$ triệu đồng. Chi phí công lao động chăm sóc $1\text{ ha}$ ngô là $15$ công, lợi nhuận thu được là $40$ triệu đồng. Biết rằng hợp tác xã chỉ huy động được tối đa $120$ công lao động. Lợi nhuận lớn nhất hợp tác xã có thể thu được là bao nhiêu triệu đồng?
\answerbox

\item Cho hệ bất phương trình bậc nhất hai ẩn có chứa tham số $m > 0$:
\[
\begin{cases}
x \ge 0,\ y \ge 0 \\
x + y \le 6 \\
y \le mx
\end{cases}
\]
Biết rằng miền nghiệm của hệ là một miền tam giác có diện tích bằng $9$. Tìm giá trị của tham số $m$.
\answerbox

\item Cho cặp số $(x; y)$ thoả mãn hệ bất phương trình:
\[
\begin{cases}
0 \le x \le 5 \\
0 \le y \le 4 \\
x + y \ge 3 \\
2x - y \le 4
\end{cases}
\]
Tìm giá trị lớn nhất của biểu thức $T = 5x - 2y$.
\answerbox
\end{enumerate}

\newpage

% ==========================================
% TRANG 4: CÂU 25 - 26 + CHỦ ĐỀ 3 (CÂU 27 - 32)
% ==========================================
\begin{enumerate}[resume]
\item Một quán giải khát pha chế hai loại đồ uống $A$ và $B$. Mỗi ly loại $A$ cần $200\text{ ml}$ sữa tươi và $50\text{ ml}$ siro hoa quả, mang lại tiền lãi $15$ nghìn đồng. Mỗi ly loại $B$ cần $100\text{ ml}$ sữa tươi và $100\text{ ml}$ siro hoa quả, mang lại tiền lãi $20$ nghìn đồng. Cửa hàng có sẵn $4\text{ lít}$ sữa tươi và $2{,}5\text{ lít}$ siro hoa quả. Hỏi số tiền lãi lớn nhất cửa hàng có thể thu được là bao nhiêu nghìn đồng?
\answerbox

\item Trong mặt phẳng toạ độ $Oxy$, cho đường thẳng bờ $d\colon 3x - 4y + 12 = 0$. Tính khoảng cách từ gốc toạ độ $O(0; 0)$ đến đường thẳng $d$.
\answerbox
\end{enumerate}

\parttitle{CHỦ ĐỀ 3. HỆ THỨC LƯỢNG TRONG TAM GIÁC VÀ ỨNG DỤNG THỰC TẾ (CÂU 27 -- 40)}

\begin{enumerate}[resume]
\item Cho tam giác $ABC$ có $AB = 5$, $AC = 8$ và góc $\widehat{BAC} = 60^\circ$. Tính độ dài cạnh $BC$.
\answerbox

\item Cho tam giác $ABC$ có cạnh $BC = 6\sqrt{3}$ và góc $\widehat{BAC} = 60^\circ$. Tính bán kính $R$ của đường tròn ngoại tiếp tam giác $ABC$.
\answerbox

\item Cho tam giác $ABC$ có độ dài ba cạnh lần lượt là $a = 13$, $b = 14$, $c = 15$. Tính bán kính $r$ của đường tròn nội tiếp tam giác $ABC$.
\answerbox

\item Cho tam giác $ABC$ vuông tại $A$ có $AB = 6$, $AC = 8$ và $BC = 10$. Tính độ dài đường trung tuyến $m_a$ kẻ từ đỉnh $A$ của tam giác $ABC$.
\answerbox

\item Để đo chiều cao của một ngọn núi đá thẳng đứng, từ hai điểm quan sát $A$ và $B$ trên mặt đất cách nhau $500\text{ m}$ cùng thẳng hàng với chân núi, người ta ngắm đỉnh núi với các góc nâng lần lượt là $30^\circ$ và $45^\circ$. Tính chiều cao ngọn núi theo đơn vị mét (làm tròn kết quả đến hàng đơn vị).
\answerbox

\item Để xác định khoảng cách giữa hai vị trí $A$ và $B$ bị ngăn cách bởi một hồ nước sâu, người ta chọn một vị trí quan sát $C$ trên bờ sao cho khoảng cách $CA = 80\text{ m}$, $CB = 120\text{ m}$ và góc nhìn $\widehat{ACB} = 120^\circ$. Tính khoảng cách $AB$ theo đơn vị mét (làm tròn kết quả đến hàng phần mười).
\answerbox
\end{enumerate}

\newpage

% ==========================================
% TRANG 5: CÂU 33 - 40
% ==========================================
\begin{enumerate}[resume]
\item Hai chiếc tàu thuỷ cùng xuất phát từ một bến cảng $O$ vào lúc $7$ giờ sáng, chạy theo hai hướng thẳng tạo với nhau một góc $60^\circ$. Tàu thứ nhất chạy với vận tốc không đổi $30\text{ km/h}$, tàu thứ hai chạy với vận tốc không đổi $40\text{ km/h}$. Hỏi sau đúng $2$ giờ kể từ lúc khởi hành, khoảng cách giữa hai tàu bằng bao nhiêu kilômét? (Làm tròn kết quả đến hàng phần mười).
\answerbox

\item Cho tam giác $ABC$ có tỉ lệ độ dài ba cạnh thoả mãn $a : b : c = 3 : 5 : 7$. Tính số đo góc lớn nhất của tam giác $ABC$ theo đơn vị độ.
\answerbox

\item Hai thôn $A$ và $B$ nằm cùng một phía của một con kênh thẳng $d$. Khoảng cách từ thôn $A$ đến bờ kênh là $AH = 3\text{ km}$, khoảng cách từ thôn $B$ đến bờ kênh là $BK = 5\text{ km}$, và khoảng cách giữa hai hình chiếu $HK = 6\text{ km}$. Người ta cần xây dựng một trạm cấp nước sạch $M$ trên bờ kênh $d$ để dẫn nước về hai thôn. Xác định chiều dài ngắn nhất của tổng đường ống dẫn nước $MA + MB$ theo đơn vị kilômét.
\answerbox

\item Cho tứ giác lồi $ABCD$ có độ dài hai đường chéo là $AC = 12$ và $BD = 16$. Biết rằng góc tạo bởi hai đường chéo bằng $30^\circ$. Tính diện tích của tứ giác $ABCD$.
\answerbox

\item Cho tam giác $ABC$ có độ dài ba cạnh lần lượt là $a = 13$, $b = 14$, $c = 15$ và bán kính đường tròn ngoại tiếp là $R$. Tính giá trị của biểu thức $P = \dfrac{a \cdot b \cdot c}{4R}$.
\answerbox

\item Một bức tranh nghệ thuật có chiều cao $3\text{ m}$ được treo trên một bức tường thẳng đứng sao cho mép dưới của bức tranh cao hơn tầm mắt của người ngắm tranh là $1\text{ m}$. Hỏi người ngắm tranh phải đứng cách bức tường một khoảng bằng bao nhiêu mét để góc nhìn bức tranh theo phương thẳng đứng đạt giá trị lớn nhất?
\answerbox

\item Cho tam giác $ABC$ có $AB = 3$, $AC = 6$ và góc $\widehat{BAC} = 120^\circ$. Tính độ dài đường phân giác trong $AD$ kẻ từ đỉnh $A$ của tam giác $ABC$.
\answerbox

\item Cho tam giác $ABC$ có độ dài ba cạnh là $a = 5$, $b = 6$ và $c = 7$. Tính bán kính $R$ của đường tròn ngoại tiếp tam giác $ABC$ (làm tròn kết quả đến hàng phần trăm).
\answerbox
\end{enumerate}

\newpage

% =========================================================================
% TRANG 6: BẢNG ĐÁP SỐ TỔNG HỢP VÀ TÓM TẮT ĐÁP SỐ CHI TIẾT
% =========================================================================
\pagestyle{dapan}
\parttitle{BẢNG ĐÁP SỐ TỔNG HỢP 40 CÂU}

\vspace{0.2cm}
\begin{center}
\begin{tabular}{|c|c||c|c||c|c||c|c|}
\hline
\textbf{Câu} & \textbf{Đáp số} & \textbf{Câu} & \textbf{Đáp số} & \textbf{Câu} & \textbf{Đáp số} & \textbf{Câu} & \textbf{Đáp số} \\
\hline
\textbf{1} & $2$ & \textbf{11} & $15$ & \textbf{21} & $14$ & \textbf{31} & $683$ \\
\hline
\textbf{2} & $3$ & \textbf{12} & $16$ & \textbf{22} & $340$ & \textbf{32} & $174{,}4$ \\
\hline
\textbf{3} & $4$ & \textbf{13} & $7$ & \textbf{23} & $1$ & \textbf{33} & $72{,}1$ \\
\hline
\textbf{4} & $2$ & \textbf{14} & $21$ & \textbf{24} & $12$ & \textbf{34} & $120^\circ$ \\
\hline
\textbf{5} & $15$ & \textbf{15} & $14$ & \textbf{25} & $550$ & \textbf{35} & $10$ \\
\hline
\textbf{6} & $45$ & \textbf{16} & $12$ & \textbf{26} & $2{,}4$ & \textbf{36} & $48$ \\
\hline
\textbf{7} & $12$ & \textbf{17} & $2600$ & \textbf{27} & $7$ & \textbf{37} & $84$ \\
\hline
\textbf{8} & $3$ & \textbf{18} & $32$ & \textbf{28} & $6$ & \textbf{38} & $2$ \\
\hline
\textbf{9} & $11$ & \textbf{19} & $10$ & \textbf{29} & $4$ & \textbf{39} & $2$ \\
\hline
\textbf{10} & $11$ & \textbf{20} & $980$ & \textbf{30} & $5$ & \textbf{40} & $3{,}57$ \\
\hline
\end{tabular}
\end{center}

\vspace{0.3cm}
\parttitle{TÓM TẮT ĐÁP SỐ CHI TIẾT CÁC BÀI TOÁN TỐI ƯU HÓA KINH TẾ}
\begin{itemize}[leftmargin=1.5em, itemsep=3.5pt]
    \item \textbf{Câu 17 (Xưởng mộc):} Sản xuất $5$ bàn trà và $3$ ghế tựa, lợi nhuận tối đa là $2600$ nghìn đồng ($2{,}6$ triệu đồng).
    \item \textbf{Câu 18 (Thuê xe):} Thuê $5$ xe loại $A$ và $4$ xe loại $B$, chi phí nhỏ nhất là $32$ triệu đồng.
    \item \textbf{Câu 20 (Dinh dưỡng):} Mua $2\text{ kg}$ thịt bò và $3\text{ kg}$ thịt lợn, số tiền ít nhất phải chi là $980$ nghìn đồng.
    \item \textbf{Câu 22 (Nông nghiệp):} Trồng $6\text{ ha}$ lúa và $4\text{ ha}$ ngô, lợi nhuận lớn nhất thu được là $340$ triệu đồng.
    \item \textbf{Câu 25 (Pha chế):} Pha chế $10$ ly loại $A$ và $20$ ly loại $B$, tiền lãi lớn nhất thu được là $550$ nghìn đồng.
    \item \textbf{Câu 35 (Trạm bơm bờ kênh):} Trạm bơm $M$ đặt tại vị trí cách hình chiếu $H$ một khoảng $HM = 2{,}25\text{ km}$, tổng chiều dài ngắn nhất là $MA + MB = 10\text{ km}$.
    \item \textbf{Câu 38 (Góc nhìn bức tranh):} Người ngắm đứng cách chân tường $2\text{ m}$ thì góc nhìn bức tranh đạt cực đại $\approx 36^\circ 52'$.
\end{itemize}

\newpage

% =========================================================================
% TRANG 7: HÌNH VẼ MINH HOẠ MIỀN NGHIỆM TIKZ TIÊU BIỂU
% =========================================================================
\parttitle{HÌNH VẼ MINH HOẠ MIỀN NGHIỆM TIKZ TIÊU BIỂU}

\vspace{0.15cm}
\noindent\textbf{1. Miền nghiệm Câu 16 -- Tứ giác diện tích $S = 12$:}
\begin{center}
\begin{tikzpicture}[scale=0.48]
    \draw[help lines, color=gray!25, step=1] (-1,-1) grid (13,13);
    \draw[-{Latex[length=2.5mm]}, line width=1.1pt] (-1.2,0) -- (13.5,0) node[below=2pt, font=\small\bfseries] {$x$};
    \draw[-{Latex[length=2.5mm]}, line width=1.1pt] (0,-1.2) -- (0,13.5) node[left=2pt, font=\small\bfseries] {$y$};
    \node[below left=2pt, font=\small\bfseries] at (0,0) {$O$};

    % Miền nghiệm tứ giác OABC
    \fill[ink!20] (0,0) -- (4,0) -- (3,3) -- (0,4) -- cycle;
    \draw[line width=1.4pt, ink] (0,0) -- (4,0) -- (3,3) -- (0,4) -- cycle;

    % Đường d1: 3x + y = 12
    \draw[line width=1.1pt, blue!80!black] (-0.3,12.9) -- (4.3,-0.9) node[pos=0.08, above right, font=\footnotesize\bfseries] {$3x + y = 12$};
    % Đường d2: x + 3y = 12
    \draw[line width=1.1pt, red!80!black] (-0.9,4.3) -- (12.9,-0.3) node[pos=0.9, above right, font=\footnotesize\bfseries] {$x + 3y = 12$};

    \fill[black] (0,0) circle (3pt);
    \fill[black] (4,0) circle (3pt) node[below=2pt, font=\footnotesize\bfseries] {$A(4;0)$};
    \fill[black] (3,3) circle (3pt) node[above right=1pt, font=\footnotesize\bfseries] {$B(3;3)$};
    \fill[black] (0,4) circle (3pt) node[left=2pt, font=\footnotesize\bfseries] {$C(0;4)$};
    \draw[dashed, gray] (3,0) node[below=2pt, font=\footnotesize] {$3$} -- (3,3) -- (0,3) node[left=2pt, font=\footnotesize] {$3$};
    \node[font=\bfseries\color{ink}] at (1.5,1.5) {$S = 12$};
\end{tikzpicture}
\end{center}

\vspace{0.25cm}
\noindent\textbf{2. Miền nghiệm Câu 17 -- Bài toán xưởng mộc (Bàn trà và ghế tựa):}
\begin{center}
\begin{tikzpicture}[scale=0.52]
    \draw[help lines, color=gray!25, step=2] (-1,-1) grid (10,19);
    \draw[-{Latex[length=2.5mm]}, line width=1.1pt] (-1.2,0) -- (10.5,0) node[below=2pt, font=\small\bfseries] {$x$};
    \draw[-{Latex[length=2.5mm]}, line width=1.1pt] (0,-1.2) -- (0,19.5) node[left=2pt, font=\small\bfseries] {$y$};
    \node[below left=2pt, font=\small\bfseries] at (0,0) {$O$};

    % Miền nghiệm OABC
    \fill[ink!20] (0,0) -- (6,0) -- (5,3) -- (0,8) -- cycle;
    \draw[line width=1.4pt, ink] (0,0) -- (6,0) -- (5,3) -- (0,8) -- cycle;

    % d1: 3x + y = 18
    \draw[line width=1.1pt, blue!80!black] (0.2,17.4) -- (6.3,-0.9) node[pos=0.08, above right, font=\footnotesize\bfseries] {$3x + y = 18$};
    % d2: x + y = 8
    \draw[line width=1.1pt, red!80!black] (-0.8,8.8) -- (8.8,-0.8) node[pos=0.9, above right, font=\footnotesize\bfseries] {$x + y = 8$};

    \fill[black] (0,0) circle (3pt);
    \fill[black] (6,0) circle (3pt) node[below=2pt, font=\footnotesize\bfseries] {$A(6;0)$};
    \fill[black] (5,3) circle (3pt) node[above right=1pt, font=\footnotesize\bfseries] {$B(5;3)$};
    \fill[black] (0,8) circle (3pt) node[left=2pt, font=\footnotesize\bfseries] {$C(0;8)$};
    \draw[dashed, gray] (5,0) node[below=2pt, font=\footnotesize] {$5$} -- (5,3) -- (0,3) node[left=2pt, font=\footnotesize] {$3$};
    \node[font=\bfseries\color{ink}] at (2.2,2.5) {Lợi nhuận max tại $B(5;3)$};
\end{tikzpicture}
\end{center}

\end{document}
